% TOP_PROBLEM: {"id": 4800018, "title": "Multiple ergodic averages — Problem 18", "category_id": 11, "category": {"id": 11, "name": "dynamical_systems", "display_name": "Dynamical Systems", "description": "Problems about long-term behavior of deterministic systems, Hamiltonian dynamics, and periodic orbits.", "slug": "dynamical-systems", "order_index": 11, "created_at": "2026-07-31T15:26:25.671Z"}, "status": "open", "problem_number": "AMR-047-0018", "set_id": 13, "statusQualification": "The source update records the two-polynomial case; the arbitrary finite family remains the target."}
% FIRST_POSTED: 2026-10-01
% ENTRY_KIND: statement_only
% UnsolvedMath ID 4800018; source catalogue code AMR-047-0018
% !TeX program = lualatex
% Definition and sources only; this file is not an LLM solution attempt.
\documentclass[11pt,a4paper]{article}
\usepackage[margin=25mm]{geometry}
\usepackage{fontspec}
\setmainfont{lmroman10-regular.otf}[BoldFont=lmroman10-bold.otf,
 ItalicFont=lmroman10-italic.otf,BoldItalicFont=lmroman10-bolditalic.otf]
\usepackage{amsmath,amssymb,mathtools}
\usepackage{unicode-math}
\setmathfont{latinmodern-math.otf}
\AtBeginDocument{\let\setminus\smallsetminus}
\usepackage{xurl,hyperref}
\hypersetup{colorlinks=true,urlcolor=blue}
\setcounter{secnumdepth}{0}
\setlength{\parindent}{0pt}
\setlength{\parskip}{5pt}
\setlength{\emergencystretch}{3em}
\begin{document}
\section*{Multiple ergodic averages — Problem 18}\label{4800018}
{\small UnsolvedMath ID 4800018 · AMR-047-0018 · Dynamical Systems · AMR Open Problem Lists}
\subsection{Definitions and mathematical statement}
Let $p_1,\ldots,p_\ell$ be integer polynomials, with
$\ell\ge1$. The family is jointly intersective if
\[
       \forall r\ge1\quad\exists n\ge1\quad
                p_i(n)\equiv0\pmod r\quad(1\le i\le\ell).
\]
The same $n$ must work for every polynomial at the
chosen modulus; requiring separate roots for each
polynomial would be weaker.

Let $(X,\mathcal B,\mu)$ be any probability space and
$T_1,\ldots,T_\ell$ any commuting invertible
measure-preserving transformations. Is joint
intersectivity sufficient to ensure that every
$A\in\mathcal B$ with $\mu(A)>0$ has
\[
           \mu\bigl(A\cap T_1^{p_1(n)}A\cap\cdots
                            \cap T_\ell^{p_\ell(n)}A\bigr)>0
                   \quad\text{for some }n\ge1?
\]
Here $T_i^kA$ is the image under the invertible
iterate, also when $k$ is negative. The positive
integer $n$ may depend on the entire system and
on $A$.

No zero-constant-term, distinctness, rational
independence, or ergodicity hypothesis is imposed.
In particular, intersective polynomials need not
have a common integer root. The requested conclusion
is positive measure of one simultaneous intersection,
without a prescribed quantitative lower bound.
The author's progress page records a result for two
polynomials; the general arbitrary-family target is
retained here. Testing periodic systems explains
the modular obstruction, but the question asks for
sufficiency in all probability-preserving systems.
\subsection{Short English statement}
Do common roots modulo every integer guarantee polynomial multiple recurrence for arbitrary commuting transformations?
\subsection{Sources}
\begin{itemize}
\item[S1] ULAM.AI and the UnsolvedMath Contributors, supplied problems.json, record 4800018 (AMR-047-0018), AMR Open Problem Lists. Catalogue identity and source transcription; CC BY 4.0.
\url{https://huggingface.co/datasets/ulamai/UnsolvedMath}.
\item[S2] Frantzikinakis - Some open problems on multiple ergodic averages (2016), Problem 18. Original source indexed by AMR-047-0018.
\url{https://arxiv.org/abs/1103.3808}.
\item[S3] Frantzikinakis, November 2025 progress update for Problem 18.
\url{https://users.math.uoc.gr/~frantzikinakis/OpenProblemsNew.html}.
\end{itemize}
\end{document}
